\section{Interpreting capacity, exposure, and generalization}
\label{sec:baselineanalysis}
The baseline validation overlay is in Appendix Figure~\ref{fig:baselineval}.

\subsection*{Evidence for overfitting, not numerical divergence}
The final validation-minus-training gaps are 0.2055, 0.5593, and 0.6360 for small, medium, and large. For the latter two, validation rises repeatedly after its minimum while training loss falls. This is consistent with overfitting, rather than loss explosion or numerical failure: the recorded losses remain finite and training continues to improve.

The gap is descriptive, not an independent estimate of generalization error: it compares fixed samples from two different splits. Its widening is more informative than its exact absolute size. No repeated-seed experiment was performed, so differences should not be labeled statistically significant.

\subsection*{Why the largest model need not win}
Medium and large have approximately 5.1 and 16.9 times small's parameter count, while receiving fewer updates and fewer repeated training targets. Large still fits the training sample strongly, suggesting that lack of representational capacity is not the immediate bottleneck. However, lower training loss is insufficient to disentangle data exposure, implicit regularization, architecture, and optimizer effects. A shared learning rate is controlled but not necessarily equally suitable for all widths/depths.

These findings favor small as the observed compute-efficient baseline. They do not demonstrate convergence of every configuration, universal superiority of small models, or a scaling law. The outcome is bounded to roughly two A10 hours per run and this music-event vocabulary.

\subsection*{Implication for the next intervention}
Checkpoint selection already avoids using deteriorated final weights: ``best within budget'' rewards an earlier successful state without pretending that all models peak simultaneously. It does not improve the best loss already reached. The next intervention therefore targets \emph{lowering the validation minimum}, not merely postponing termination. Increasing dropout provides a controlled test of whether larger models can use their capacity with less reliance on training-specific feature combinations. Validation remains the adaptation signal; the test split is not consulted during this decision.
